\section{Task 3: Dataset, Full Experimental Log and Limitations}
\label{app:t3}
\renewcommand{\expprefix}{C}

\subsection{Dataset}

Table~\ref{tab:t3-data} summarises the data.

\begin{table}[h]
\caption{Task 3 dataset statistics.}
\label{tab:t3-data}
\centering
\begin{tabular}{ll}
\hline
Quantity & Value \\
\hline
Training / test flows & 1{,}285 / 327 \\
Classes & 10 (5 applications $\times$ voice/video) \\
Class sizes & 40 (GoogleMeet\_voice) to 256 (Discord\_video) \\
Raw fields per flow & 5 packet lengths, 5 relative timestamps \\
Packet lengths & 26--1{,}242 bytes \\
Five-packet time span & 48\,$\mu$s to 23.21\,s (median 571\,$\mu$s) \\
Source calls (training) & $\approx$400 (estimated, see Exp.~2; identifiers not released) \\
Train/test adversarial-validation AUC & 0.529 (exchangeable) \\
Constant-size flows & 125 (Discord: 58 voice, 37 video) \\
Audio-only windows (all five packets $<300$\,B) & 726 flows (56.5\%) \\
Zoom audio-only windows & 160 (80 voice, 80 video) \\
Official metrics & accuracy, precision, recall, F$_\beta$ (weighted) \\
Development metric & macro-F1 \\
\hline
\end{tabular}
\end{table}

\subsection{Organisation of the Log}

Twenty-five numbered experiments were run in a logical sequence; we group them
by theme. Within this appendix, ``Experiment~$n$'' refers to Experiment~C$n$. Unless noted, scores for Branch~B are grouped cross-validation scores
(Experiment~2). Branch~A was developed earlier, under stratified 5-fold
cross-validation with nested tuning, and its numbers are labelled accordingly
(see Section~\ref{sec:t3-limits}). Metrics are accuracy, macro-F1 and, for
Branch~A, macro-recall.

\subsection{Data and validation protocol}

 
\begin{explog}{1}{Exploratory analysis of the flow records}
\item[Hypothesis] Before any modelling, structural properties of the 5-packet records
(duplicates, train/test shift, label determinism, protocol markers) decide how
much is learnable and how evaluation must be built.
\item[Method] Integrity checks; per-column Kolmogorov--Smirnov tests and adversarial
validation (a RandomForest separating train rows from test rows); per-class
packet-length and inter-arrival analysis against protocol reference bands;
mutual information and ANOVA-$F$ per engineered feature; embeddings coloured
separately by application and by call mode; flat versus factorised
(application $\times$ mode) classifiers.
\item[Result] 1{,}285 training and 327 test flows, no missing values,
\texttt{relative\_time\_0} identically zero. Adversarial-validation AUC 0.529.
Only 2 of 1{,}136 distinct length tuples map to more than one class, but they
cover 95 rows (7.4\%). Packets in the 1100--1600\,B band make up 10.7\% of
video-class packets and 0.0\% of voice-class packets. Factorised, the
application is recoverable at 0.912 and the call mode only at 0.888. Flows per
call extrapolated to 100 held-out calls predict $\approx$321 test flows against
327 supplied.
\item[Inference] Train and test are exchangeable, so cross-validation is a meaningful
proxy. Call mode, not application, is the harder factor. Because most length
tuples are unique, tuple-to-label determinism says little about held-out rows
until it is tested on them (Experiment~13).
\item[Verdict] Kept as the standing diagnostic; motivated the feature bands and the
validation design.
\end{explog}

\begin{explog}{2}{Plain versus grouped validation, and negative controls}
\item[Hypothesis] The 1{,}285 training flows come from an estimated 400 source calls and no call
identifier is released, so near-identical sibling flows will straddle the
train/validation boundary under ordinary stratified folds and may inflate scores.
If folds and cached features are leak-free, shuffled labels must score near
chance and a prior-only predictor must match the majority share.
\item[Method] Two schemes were run side by side for every Branch~B model: \emph{plain}
(repeated stratified 5-fold, 3 repeats) and \emph{grouped}
(\texttt{StratifiedGroupKFold}, 5 splits, 2 shuffled repeats $=10$ folds) over a
pseudo-group key equal to the packet-length tuple rounded to 8-byte granularity
joined with $\mathrm{round}(\log_{10}\mathrm{duration},1)$. Fold indices were
stored. Controls: N1, grouped CV with permuted labels; N2, prior-only predictor.
\item[Result] The key yields 1{,}107 pseudo-groups; 253 flows (20\%) share a group with
at least one other flow. Plain and grouped scores of matched runs differ by up
to about two points, and not always in the same direction (e.g.\ Experiments~10
and~14). N1 scores 0.1549 (expected band 0.08--0.20); N2 scores 0.1992
(expected 0.1992).
\item[Inference] Sibling leakage is moderate at this data size but cannot be ruled out per
model, so grouped scores are the ones every selection decision relies on. Folds
and feature cache are leak-free. Two caveats. First, the pseudo-key is
conservative in one direction only: 1{,}107 pseudo-groups against an estimated
400 calls means the key splits some calls, so grouped CV may still leak siblings
(Experiment~15 finds 1{,}024 singleton clusters). Second, for TabICL at full
width the grouped score (0.8389) is \emph{higher} than the plain score (0.8195),
the opposite of what sibling leakage predicts; we have not isolated the cause and
treat the two protocols as differing by seed and fold structure as much as by
leakage.
\item[Verdict] Grouped CV adopted as the decisive estimate; both controls pass.
\end{explog}

\subsection{Baselines, model search and features}

 
\begin{explog}{3}{Untuned baseline and large-scale model search}
\item[Hypothesis] A broad search across model families finds a configuration that clearly
outperforms an untuned gradient-boosted baseline.
\item[Method] Untuned LightGBM on the engineered features (plain 5-fold) as control.
Then 2{,}377 points scored (XGBoost 962, LightGBM 626, RandomForest 588, CatBoost
61, ExtraTrees 36, SVM 27, MLP 23, $k$-NN 22, HistGradientBoosting 20, logistic
regression 10, PacketFormer 2) in 83.6 container-hours, one container per grid
point, then Optuna~\cite{optuna} (TPE, ask/tell batching) around the grid
winners. Class weighting was itself a searched parameter. Grouped CV.
\item[Result] Untuned LightGBM 0.810 accuracy (plain protocol, so not directly comparable
with the grouped scores below). Optuna: LightGBM 180 trials $\rightarrow$ 0.8249,
XGBoost 180 $\rightarrow$ 0.8202, RandomForest 150 $\rightarrow$ 0.8023; the
CatBoost~\cite{catboost} study stopped when a single trial exceeded the 3{,}600\,s
limit. Table~\ref{tab:t3-leader} gives the best point per family. The winning
LightGBM point uses 800 trees, 31 leaves, learning rate 0.12, subsample 0.8,
column sample 0.8, minimum child samples 10 and no class weighting.
\item[Inference] Boosted trees lead; distance, kernel, linear and neural models trail by
9--11 points. The untuned baseline was scored under the plain protocol and the searched points
under grouped CV, so the gain attributable to the search is not isolated; the
grouped-CV gap between the tuned LightGBM (0.8249) and the default-setting
LightGBM of Experiments~4 and~17 (0.8152, 0.8156) is about one point.
% TODO: run untuned LightGBM on the stored grouped folds and replace this sentence with the matched number
\item[Verdict] LightGBM (0.8249 grouped, macro-F1 0.8069) kept as the control for all
later Branch~B comparisons. The ``LightGBM control'' appears with four scores in
this section (0.8249 here; 0.8152 in Experiment~4; 0.8210 in Experiment~14;
0.8156 in Experiment~17) because they come from different configurations and
repeat counts; Table~\ref{tab:t3-control} lists them.
% TODO: fill the config column from the logs
\end{explog}

\begin{table}[t]
\caption{Best point per family, grouped CV (Experiment~3).}
\label{tab:t3-leader}
\centering
\begin{tabular}{rlcc}
\hline
Rank & Model & Grouped accuracy & Macro-F1 \\
\hline
1 & LightGBM & 0.8249 & 0.8069 \\
2 & XGBoost & 0.8218 & 0.8067 \\
3 & HistGradientBoosting & 0.8156 & 0.7959 \\
4 & ExtraTrees & 0.8125 & 0.8009 \\
5 & CatBoost & 0.8070 & 0.7895 \\
6 & RandomForest & 0.8023 & 0.7866 \\
7 & PacketFormer (pretrained) & 0.7346 & 0.7262 \\
8 & PacketFormer (scratch) & 0.7331 & 0.7239 \\
9 & MLP & 0.7292 & 0.6792 \\
10 & SVM (RBF) & 0.7268 & 0.6595 \\
11 & Logistic regression & 0.7144 & 0.6666 \\
12 & $k$-NN & 0.7132 & 0.6747 \\
-- & Majority class & 0.1992 & 0.0332 \\
\hline
\end{tabular}
\end{table}

 \begin{table}[t]
\caption{The four LightGBM scores used as ``control'' in this section.}
\label{tab:t3-control}
\centering
\begin{tabular}{llc}
\hline
Experiment & Configuration & Grouped accuracy \\
\hline
3 & Optuna best point, 241 columns & 0.8249 \\
4 & matched full-width run, 241 columns & 0.8152 \\
14 & best point used for calibration & 0.8210 \\
17 & standing configuration, three repeats & 0.8156 \\
\hline
\end{tabular}
\end{table}
\begin{explog}{4}{Feature width and label-aware blocks}
\item[Hypothesis] Pruning the 241 deterministic columns by mutual information should remove
noise at 1{,}285 rows; class-conditional Markov, density and neighbourhood
features computed inside the CV loop should add signal beyond the label-free
blocks.
\item[Method] Matched LightGBM and XGBoost runs on the full 241 columns, top-120 and
top-60 columns, full repeat budget; then matched runs with the label-aware blocks
switched on and off. % TODO: state whether the MI ranking was computed inside each fold
\item[Result] LightGBM grouped accuracy 0.8152 (full), 0.8097 (top-120), 0.7895 (top-60);
macro-F1 0.7974, 0.7918, 0.7706. XGBoost 0.8016 (full) against 0.7907 and 0.7911.
Label-aware blocks: LightGBM 0.8152 $\rightarrow$ 0.7969, XGBoost 0.8016
$\rightarrow$ 0.7883, at every width tested.
\item[Inference] The full set beats both subsets in every matched pair, so pruning discards
usable signal. At 1{,}285 rows the label-aware blocks contribute variance rather
than signal (about 1.8 accuracy points on LightGBM) even after the inner
out-of-fold fix.
\item[Verdict] Full 241-column label-free set kept for Branch~B; label-aware blocks
dropped.
\end{explog}

\begin{explog}{5}{Compact 31-feature set and two added features (Branch A)}
\item[Hypothesis] Branch~A needs only a compact representation (raw lengths and times,
differences, order statistics, shape descriptors); single features found by
symbolic search or by inspecting misclassified Discord and Google Meet video flows
may still add signal.
\item[Method] 29 base columns (Stage~0 of Section~\ref{sec:t3-arch}) plus two features added one at a time:
$\log$ of the fourth packet length (found by genetic-programming symbolic
regression), and the position of the largest packet-to-packet size increase
(separating an instant jump to a plateau from a gradual ramp, found by manual
inspection of video misclassifications).
% TODO: state the protocol and metric of the two gains (stratified 5-fold nested? macro-recall or macro-F1?)
\item[Result] An isolated gain in the validated score of $+0.0017$ for the log feature and
$+0.0014$ for the jump position. Total 31 features.
\item[Inference] Small but individually reproducible gains; the jump position encodes the
temporal shape that static order statistics miss.
\item[Verdict] Both kept; Branch~A uses 31 features.
\end{explog}

\subsection{Branch A: tree ensemble}

 
\begin{explog}{6}{Probability correction on the averaged ensemble}
\item[Hypothesis] Averaged tree probabilities are over-confident and carry the training class
prior; softening them with a temperature and dividing out the prior should help a
macro-averaged criterion.
\item[Method] Out-of-fold probabilities from the three models (Stage~1 settings) under
stratified 5-fold (seed 42); equal-weight average; temperature $T$ on the
log-probabilities; division by the empirical class prior; argmax. $T$ was fixed at
2.0 after the temperature studies of this branch. % TODO: add the T sweep and the with/without-prior ablation as a small table
\item[Result] Stage-0 macro-recall 0.8286.
\item[Inference] Reference point for the decision layer of Experiment~7.
\item[Verdict] Kept.
\end{explog}

\begin{explog}{7}{Per-class decision bias and margin-triggered specialist}
\item[Hypothesis] Four classes (Zoom\_video, GoogleMeet\_video, Zoom\_voice, Discord\_video)
were identified as weak, and the residual confusions among five classes
(GoogleMeet voice/video, Discord voice/video, Messenger\_video) are genuine
two-way ties with a sample-dependent partner. An additive log-probability offset
on the weak classes, plus a specialist model for tied cases, should recover them.
Zoom\_video is the class sacrificed, so leaving its offset unconstrained should be
harmful.
\item[Method] Within five outer folds, offsets for the four weak classes were tuned by
coordinate ascent on the tuning rows only (coarse grid $-1.5$ to $1.5$ in 31 steps,
then a $\pm0.15$ fine grid, up to 10 passes). A specialist LightGBM trained only on
the five confusable classes replaced the prediction whenever the top-1 and top-2
classes were both confusable and the probability margin was below
$\tau\in\{0,0.05,\dots,0.25\}$; $\tau$ and the specialist were fitted in an inner
four-fold loop. Final constants are medians across outer folds. The Zoom\_video
offset was also capped at $\geq-0.10$ and compared with unrestricted optimisation.
\item[Result] Macro-recall 0.8286 $\rightarrow$ 0.8379 (nested). Final offsets:
Zoom\_video $-0.10$, GoogleMeet\_video $+0.20$, Zoom\_voice $+0.12$,
Discord\_video $+0.38$; $\tau=0.05$. Uncapped Zoom\_video offset: 0.8358; capped:
0.8379.
\item[Inference] A gain of 0.0093 over the corrected ensemble, obtained by shifting the
decision boundary toward recall (the precision cost is measured in Experiment~8).
Unconstrained optimisation pushes Zoom\_video further down for a worse net result.
% TODO: add the confusion matrix that identified the 4 weak / 5 confusable classes, and offsets-only vs specialist-only arms
\item[Verdict] Kept as the stand-alone Branch~A configuration, with the cap.
\end{explog}

\begin{explog}{8}{F1 view of Branch A}
\item[Hypothesis] Recall-oriented offsets improve macro-recall at some cost in precision;
macro-F1, the scored metric, will show where.
\item[Method] Macro-, micro- and weighted-F1 and per-class recall and F1 on the nested
predictions of Experiment~7.
\item[Result] Macro-F1 0.8033, micro-F1 0.8078, weighted-F1 0.8131, against macro-recall
0.8379. Zoom\_voice recall 0.800 but F1 0.549; GoogleMeet\_voice recall 0.950, F1
0.724; Zoom\_video recall 0.552, F1 0.691; WhatsApp\_voice recall 0.988, F1 0.975.
\item[Inference] Zoom\_voice precision is about 0.42: the offsets buy recall with false
positives, so the macro-recall figure overstates the criterion actually scored.
\item[Verdict] The branch is reported on both metrics; its outputs enter the blend as
corrected probabilities \emph{without} the offsets or the specialist
(Experiment~25).
\end{explog}

\begin{explog}{9}{Single-thread determinism}
\item[Hypothesis] Parallel fitting should not change which hyperparameter pair the inner search
selects.
\item[Method] Estimators were run with \texttt{n\_jobs=-1} and \texttt{n\_jobs=1}; the
selected (temperature, margin-threshold) pair and the final score were compared.
\item[Result] Parallel execution introduced floating-point non-determinism that shifted
the selected pair and moved the final score by up to about 0.01.
\item[Inference] Differences of this size are as large as several of the gains reported in
this section.
\item[Verdict] All estimators run with \texttt{n\_jobs=1}.
\end{explog}

\subsection{Branch B: tabular in-context learning}

 
\begin{explog}{10}{TabICL against the tree control}
\item[Hypothesis] A pretrained in-context tabular classifier is competitive on 1{,}285 rows
without a hyperparameter search.
\item[Method] TabICL at default settings on the full 241 columns, top-120 and top-60,
against the LightGBM control of Experiment~3; paired McNemar test~\cite{mcnemar}.
TabPFN~v2~\cite{tabpfn} could not be run initially because its weights require
licence acceptance through an interactive login; it was later run on the published
checkpoint.
\item[Result] Full width 0.8389 grouped (plain 0.8195, macro-F1 0.8232, 83\,s); top-120
0.8280 (0.8117); top-60 0.8272 (0.8117); LightGBM control 0.8249 (0.8069). Against
the control TabICL wins 66 rows and loses 49 ($p=0.135$). TabPFN~v2: 0.8311
(top-120), 0.8230 (full), 0.8191 (top-60), degrading with width.
\item[Inference] All three TabICL widths exceed the best of 2{,}377 searched points with
three arms and no tuning, but the paired difference is not established at this
sample size. TabICL is also the best foundation model tried.
\item[Verdict] Adopted as the second branch on the strength of the point estimate and
its consistency across widths.
\end{explog}

\begin{explog}{11}{TabICL parameters and seed bagging}
\item[Hypothesis] The internal ensemble size, the softmax temperature and bagging over seeds
can add accuracy cheaply.
\item[Method] Eighteen arms over \texttt{n\_estimators}, \texttt{softmax\_temperature}
and bag count, full repeat budget.
\item[Result] \texttt{n\_estimators}=16: 0.8405 grouped (macro-F1 0.8268), $+0.0016$ over
the default 8 against a noise floor of 0.0031; 32: 0.8358. The softmax temperature
(0.75/0.90/1.05) gives identical accuracy to four decimals. Five seed bags of 8
estimators: 0.8350 (sd 0.0016) against 0.8389 unbagged (sd 0.0031).
\item[Inference] The internal ensemble has an interior optimum; temperature cancels out of the
argmax under logit averaging and acts as a calibration parameter only; bagging
halves the variance but smooths toward a slightly worse consensus because the model
already ensembles internally.
\item[Verdict] None of the three changes is established on accuracy grounds; the
runs that feed the blend nevertheless use 16 estimators.
\end{explog}

\subsection{Routes that did not help}

 
\begin{explog}{12}{Auxiliary corpus: pretraining, probe features and direct transfer}
\item[Hypothesis] A small transformer over the five packets, pretrained on a larger public
capture corpus that contains all five applications, transfers representations that
the trees cannot learn from 1{,}285 rows; frozen-encoder probes add information
absent from the engineered features; and a classifier trained purely on the corpus
recognises the application on competition flows if the two corpora are
commensurable.
\item[Method] The corpus (MIRAGE-AppAct-2024~\cite{mirage}, 6.97\,GB, 20 applications) was reduced to
payload lengths, gaps and direction for the first 20 packets: 14{,}975 UDP flows,
or 90{,}610 flows with TCP included (the larger set was adopted, with a transport
flag). \emph{Pretraining:} PacketFormer tokenises each packet (eight continuous
descriptors, a 24-bin size embedding, positional embedding, \texttt{[CLS]}); 3
pre-norm layers, 4 heads, $d_\mathrm{model}=96$; 8{,}000 steps, batch 256, 347\,s on
one GPU; masked-packet modelling, a contrastive objective over two traffic-plausible
augmentations, a domain-adversarial head against unlabelled competition flows and
random truncation to 5/8/12/20 packets. No competition label was used.
\emph{Probes:} 46 probe columns (application family, a video-likeness proxy,
bitrate and pacing with residuals, a transport fingerprint, PCA components) appended
to the full feature set; matched runs with the block on and off, for LightGBM,
XGBoost, RandomForest and TabICL at 8, 16 and 32 estimators. \emph{Transfer:} a
gradient-boosted model trained on 7{,}655 UDP media flows of the five applications
was applied to the competition training set, with lengths as stored and with 28
bytes of IP+UDP header added; the corpus metadata was inspected for an activity
label.
\item[Result] Pretraining: from scratch 0.7331 grouped accuracy (macro-F1 0.7239),
pretrained 0.7346 (0.7262), a gain of 0.0015; final pretraining losses: total
2.664, masked-bin 1.146, masked-gap 0.959, contrastive 2.074, domain 0.016. Probes
(grouped accuracy): LightGBM 0.8144 $\rightarrow$ 0.8144, XGBoost 0.8082
$\rightarrow$ 0.7996, RandomForest 0.7778 $\rightarrow$ 0.7782; TabICL with probes
0.8257 / 0.8288 / 0.8300 against 0.8389 (8) and 0.8405 (16) without, best arm
against best arm $-0.0105$; the transport-fingerprint probe had a degenerate target
and emitted no features. Transfer: application accuracy 0.3183 and 0.2988
respectively, against 0.3844 for the majority class; the published archive carries
only the package name and byte/packet counters, so no activity manifest exists.
\item[Inference] Five tokens give attention little to work with and 1{,}285 labels give
fine-tuning little to learn from. The probe block is neutral-to-harmful for three
model families and marginal for a fourth. Transfer falls below the majority
baseline: the domain gap at the level of packet sizes is total, which also explains
the null probe result.
\item[Verdict] Auxiliary-data route closed. PacketFormer kept only as a diversity member
of the model pool; probe block retired.
\end{explog}

\begin{explog}{13}{Exact and tolerance-based length-tuple lookup}
\item[Hypothesis] Because the label is almost determined by the length tuple (Experiment~1), a
lookup table over training tuples can override the model on rows it matches.
\item[Method] Coverage and purity of tuple matches at exact, $\pm1$, $\pm2$ and $\pm4$
byte tolerance (gate: more than 5\% coverage at more than 0.95 accuracy), followed
by the differential that decides the question: model accuracy against lookup
accuracy on the covered rows.
\item[Result] Exact matching covers 7.6\% of test rows (5.5\% of CV rows) at 0.972 accuracy
and passes the gate. On the 71 covered CV rows the model is already at 1.0000
against the lookup's 0.9718; overall accuracy with the override changes by
$-0.0016$. At $\pm4$ bytes: 232 rows, model 0.9181, lookup 0.8922, change
$-0.0047$.
\item[Inference] The model is already perfect on every exactly matched row; a lookup can only
take rows from a model that answers them correctly. The near-determinism seen in
Experiment~1 is a property of unique tuples, not of exploitable structure.
\item[Verdict] Dropped.
\end{explog}

\begin{explog}{14}{Calibration and decision-layer adjustments}
\item[Hypothesis] Confidence can be repaired without moving accuracy; a per-class logit bias,
or deferring low-margin rows, lifts accuracy, particularly inside the Zoom pair.
\item[Method] Temperature~\cite{guo}, vector and Dirichlet~\cite{kull} calibration on the
grouped out-of-fold predictions of the best LightGBM point; the selection rule was
amended mid-round to the best negative log-likelihood among calibrators within one
point of raw accuracy, because the first rule selected Dirichlet and silently gave
up accuracy. Nested logit-bias fits on raw and on calibrated probabilities, and
margin-based deferral at 0.10 and 0.35.
\item[Result] Raw: accuracy 0.8210, NLL 0.7931, ECE 0.1012, mean confidence when wrong
0.7833. Temperature: 0.8210, 0.5628, 0.0451, 0.6095. Vector: accuracy 0.8031, NLL
0.5898, ECE 0.0359. Dirichlet: 0.8101, 0.5450, 0.0349. Decision layer: raw argmax
0.8210 (Zoom-pair accuracy 0.6548); bias on raw 0.8226 (5 rows changed); bias on
calibrated 0.8187 (18 rows); deferral 0.8202 at 0.10 and 0.8117 at 0.35. The fitted
bias vector is small (max $|b|=0.10$) and does not push toward Zoom\_video.
\item[Inference] Temperature scaling cuts NLL by 29\% and ECE by 55\% at no accuracy cost; the
richer calibrators buy slightly more calibration for 1--2 accuracy points. The best
decision-layer arm moves five rows for $+0.0016$, inside the noise floor; a global
per-class prior shift cannot recover the Zoom-pair gap, because buying Zoom rows
costs rows elsewhere.
\item[Verdict] Temperature scaling kept; decision-layer adjustments dropped.
\end{explog}

\begin{explog}{15}{Recovering call structure from the flows}
\item[Hypothesis] If flows from one call can be grouped, aggregating predictions within a group
should resolve ambiguous flows from their siblings.
\item[Method] A same-call scorer, clustering of its output, and aggregation of per-flow
predictions within clusters by log-probability sum or confidence-weighted vote,
plus an oracle run using the pseudo-groups as if they were true calls.
\item[Result] Same-call ROC-AUC 0.998 on held-out groups; cluster application purity 0.997,
label purity 0.989, adjusted Rand index 0.775 against the pseudo-groups; 1{,}024 of
1{,}107 clusters are singletons, the largest holds 13 flows. Aggregation: 0.8327
(log-probability sum, $-0.0054$) and 0.8304 (weighted vote, $-0.0078$) against
0.8381; oracle grouping 0.8311 ($-0.0070$).
\item[Inference] Call structure is easy to recover and worthless: even with perfect knowledge
of which flows share a call, aggregation lowers accuracy, because siblings are rare
and where they exist they already agree with the model.
\item[Verdict] Branch closed on a measured ceiling.
\end{explog}

\begin{explog}{16}{Combining models inside Branch B, and selecting the Branch B model}
\item[Hypothesis] Combining diverse learners beats the best single model; a weighted
combination of TabICL and a tree model, or a larger pool, beats the pre-specified
TabICL pair.
\item[Method] Hill-climbed soft-vote weights (also scored nested: fitted on four fifths of
the out-of-fold rows, scored on the fifth), stacking with a logistic meta-learner
(inner CV) and rank averaging, in two rounds: the gradient-boosted pool, then the
pool with TabICL. Then 17 candidate blends over TabICL and the five other families
on the same grouped folds with and without the Zoom rule, with paired bootstrap and
McNemar comparisons against the earlier standing model (TabICL, flat view only).
The pre-specified pair (flat and hierarchical views averaged) was kept as the
default, because the best-scoring blend is the maximum of seventeen candidates
scored on the same 1{,}285 rows.
\item[Result] Round one: best single LightGBM 0.8210; nested hill-climb 0.8195; stacking
0.8187; rank averaging 0.7735; the fitted hill-climb score 0.8319 is 1.2 points
above its nested score. Round two: single TabICL 0.8381; stacking 0.8327; nested
hill-climb (TabICL $\times$0.75, LightGBM $\times$0.25) 0.8311; fitted 0.8451; rank
averaging 0.7837. Candidate blends, accuracy / macro-F1 without the rule, and
macro-F1 with it: TabICL $\times$2 + LightGBM (flat+hier) 0.8381 / 0.8216, 0.8319;
\textbf{TabICL flat+hier average 0.8358 / 0.8201, 0.8292}; TabICL flat only 0.8381
/ 0.8262, 0.8274; TabICL + LightGBM + XGBoost 0.8342 / 0.8147, 0.8231; all six
families 0.8272 / 0.8082, 0.8172. Shipped against standing: $+0.0030$ macro-F1,
interval $[-0.0130,+0.0201]$, $P(\Delta>0)=0.630$, McNemar 46/51 ($p=0.685$);
observed maximum against standing: $+0.0058$, $[-0.0114,+0.0231]$, 0.744, 53/54
($p=1.000$).
% TODO: the flat-only TabICL row (0.8381/0.8262) differs from Exp. 10/11 (0.8389/0.8232, 0.8405/0.8268); add one line saying these are different rounds/seeds
\item[Inference] At 1{,}285 rows a fitted blend weight is a memorised weight: only nested
estimates are admissible for choosing. Neither paired difference in model
selection is established; the durable result of the round is the diagnosis of which
flows are unwinnable (Experiment~18), not the choice among near-equal blends.
\item[Verdict] TabICL flat and hierarchical views averaged, with the Zoom rule (grouped
accuracy 0.8342, macro-F1 0.8292), is Branch~B.
\end{explog}

\subsection{Error geometry and the Zoom rule}

 
\begin{explog}{17}{Where the errors are}
\item[Hypothesis] The residual errors concentrate in a small number of identifiable confusions
rather than being spread evenly.
\item[Method] Three-repeat grouped CV of the standing gradient-boosted configuration;
errors split by label dimension and by packet-size regime (whether all five packets
lie in the audio band, below 300 bytes).
\item[Result] Accuracy 0.8156, macro-F1 0.7982, 237 errors. Application accuracy 0.9035;
mode accuracy 0.8934 (0.9027 given the application is right). Audio-only windows:
726 flows (56.5\%) at 0.7590 with 175 errors; windows with a larger packet: 559
flows at 0.8891 with 62 errors. Zoom\_voice against Zoom\_video alone accounts for
68 of the 237 errors; Zoom\_voice recall is 0.34.
\item[Inference] Errors concentrate in audio-only windows and in one pair.
\item[Verdict] Redirected the work from model search to the Zoom pair.
\end{explog}

\begin{explog}{18}{Separability of voice and video inside the audio-only window, and the score ceiling}
\item[Hypothesis] A video call whose first five packets are all audio-sized produces a record
indistinguishable from a voice-call record; whether the pair is separable depends on
the application.
\item[Method] Per application, over flows whose window is audio-only, grouped-CV AUC of
video against voice (gradient boosting over all 241 features); threshold sweep of
the audio band from 200 to 500 bytes; upper bounds on accuracy and macro-F1 assuming
every other flow is classified perfectly.
\item[Result] Table~\ref{tab:t3-ambig}. For Zoom the subset is exactly balanced (80 voice,
80 video) and the AUC of 0.551 is not distinguishable from a coin flip at $n=160$.
Every threshold from 250 to 500 bytes isolates the same 80/80 subset (75/76 at 200
bytes). Upper bounds: accuracy 0.9377 (80 of 1{,}285 flows unwinnable) and macro-F1
0.9364, reached by sending the whole subset to Zoom\_voice (Zoom\_voice F1 0.667,
Zoom\_video F1 0.697). A macro-F1 of 0.90 would need the other eight classes to
average F1 $\geq0.954$; they average 0.868.
\item[Inference] The Zoom pair is information-limited by the five-packet window, not by the
model; the 300-byte boundary is a property of the data, not of a tuned constant. This
is the Task~3 counterpart of Task~1's perfect-rescorer ceiling (Exp.~A15): the headline metric cannot go past 0.9364, so the remaining
distance is in the other classes.
\item[Verdict] A macro-F1 of 0.90 is not attainable on this dataset; the quantity of
interest becomes how much of the structural maximum is captured.
\end{explog}

\begin{table}[t]
\caption{Audio-only windows by application (Experiment~18). The four applications
shown account for 646 of the 726 audio-only flows.}
% TODO: add the WhatsApp row (the remaining 80 flows) once confirmed from the logs
\label{tab:t3-ambig}
\centering
\begin{tabular}{lccc}
\hline
Application & Audio-only flows & of which video & Grouped AUC (video vs voice) \\
\hline
Discord & 290 & 52 & 0.880 \\
Google Meet & 85 & 45 & 0.826 \\
Messenger & 111 & 19 & 0.879 \\
Zoom & 160 & 80 & 0.551 \\
\hline
\end{tabular}
\end{table}
 
\begin{explog}{19}{Derivation and measurement of the Zoom rule}
\item[Hypothesis] On an evenly split, inseparable subset every assignment has the same accuracy
but not the same macro-F1; sending the whole subset to the smaller class recovers its
recall at no cost to the larger one.
\item[Method] Derivation: with all other flows correct, sending the 160 audio-only Zoom
flows to Zoom\_voice gives F1 0.667 (voice) and 0.697 (video; recall $92/172$),
summing to 1.364, whereas sending them to Zoom\_video gives 0 and 0.811, summing to
0.811. Rule: if a flow is predicted to be Zoom (either class) and all five packets
are below 300 bytes, label it Zoom\_voice. The rule has no fitted parameter: its
direction follows from the 80/80 balance and the class sizes, and its threshold is
the audio band already defined from the protocol analysis. Out-of-fold evaluation on
the flat, hierarchical and averaged views of the round-three base learner;
% TODO: name the base learner used for these numbers
bootstrap over 2{,}000 resamples of the averaged model.
\item[Result] Macro-F1 0.7942 $\rightarrow$ 0.8080 (flat), 0.7978 $\rightarrow$ 0.8092
(hierarchical), 0.8041 $\rightarrow$ 0.8181 (average). Bootstrap gain on the
average: $+0.0140$, 95\% interval $[0.000,+0.028]$, $P(\mathrm{gain}>0)=0.975$;
accuracy 0.8233 $\rightarrow$ 0.8241; Zoom\_voice F1 0.384 $\rightarrow$ 0.580.
\item[Inference] The gain comes from Zoom\_voice recall, with accuracy essentially unchanged, as
the derivation predicts. The bootstrap interval touches zero at its lower end. The
rule was derived for macro-F1; on accuracy it is neutral (0.8233 $\rightarrow$
0.8241), so it should not be expected to move an accuracy-type score.
\item[Verdict] Kept as a deterministic post-processing step.
\end{explog}

\begin{explog}{20}{Does the rule hold across model families?}
\item[Hypothesis] A rule derived from a class-balance measurement, not fitted, should help every
hypothesis class.
\item[Method] Six base learners, each in three views (flat ten-class, hierarchical,
fifteen-class mixture in which video is split by window type), grouped folds, with
and without the rule.
\item[Result] Table~\ref{tab:t3-rule}: all six families improve, in every view. For TabICL the matched flat+hierarchical average goes from 0.8201 to 0.8292 ($+0.0091$); the $+0.0030$ implied by the table's best-view columns understates it. % TODO: recompute all six rows with the same view in both columns
\item[Inference] The rule reflects the data rather than the estimator.
\item[Verdict] Kept.
\end{explog}

\begin{table}[t]
\caption{Macro-F1 per family without and with the Zoom rule (Experiment~20); each
column reports the best of the three views, chosen separately per column, so the
two columns are not a matched-view comparison.}
\label{tab:t3-rule}
\centering
\begin{tabular}{lcc}
\hline
Base learner & Without rule & With rule \\
\hline
TabICL & 0.8262 & 0.8292 \\
LightGBM & 0.8066 & 0.8169 \\
XGBoost & 0.7965 & 0.8044 \\
ExtraTrees & 0.7815 & 0.7885 \\
CatBoost & 0.7735 & 0.7804 \\
RandomForest & 0.7640 & 0.7725 \\
\hline
\end{tabular}
\end{table}
 
\begin{explog}{21}{Fitted corrections against the derived rule}
\item[Hypothesis] Hierarchical heads, a finer label view, balanced weights or a fitted decision
layer add to the rule.
\item[Method] Each change compared on grouped out-of-fold predictions; fitted parameters
scored nested.
\item[Result] Hierarchical application $\times$ mode heads: $+0.004$ macro-F1 alone and a
useful second view. Fifteen-class mixture view: 0.8063 alone, above the flat view.
Average of flat, hierarchical and mixture: 0.8085, 0.8185 with the rule. A single
global mode head in place of per-application heads: AUC 0.9661 against 0.9680, no
difference. Class-balanced sample weights in every head: flat $+0.005$, ensemble
$-0.003$. Per-class log-probability bias fitted for macro-F1: in sample 0.8124,
nested 0.7956 against 0.7982 raw. A single global mode-bias parameter: nested 0.8173
against 0.8181 unbiased.
\item[Inference] Both fitted decision-layer variants gain in sample and lose nested; at 1{,}285
rows a fitted threshold is a memorised threshold, which is why the derived rule
survives where fitted corrections do not.
\item[Verdict] Hierarchical and mixture views kept as ensemble members; weights and
fitted biases dropped.
\end{explog}

\subsection{Combining the two branches}

 
\begin{explog}{22}{Paired out-of-fold diagnostics}
\item[Hypothesis] The two branches, built on different features (31 against 241 columns) and
different learner families, make partly different errors, so blending can help.
\item[Method] Raw Branch~A probabilities (temperature $T=2$, prior-corrected) against the
averaged Branch~B flat and hierarchical probabilities on the same 1{,}285 training
rows.
\item[Result] Macro-F1 / accuracy: Branch~A 0.7900 / 0.7961, Branch~B 0.8201 / 0.8358. The
two agree on 84.98\% of rows; both are correct on 75.33\% and both wrong on 12.14\%;
only Branch~A is correct on 4.28\% (55 rows) and only Branch~B on 8.25\% (106
rows). Branch~A's 55 unique wins are led by Zoom\_voice (27), GoogleMeet\_voice (8)
and GoogleMeet\_video (7); Branch~B's 106 by Discord\_video (30), Discord\_voice (28)
and Zoom\_video (27). Zoom\_voice recall is 0.712 for Branch~A against 0.438 for
Branch~B; Zoom\_video recall is 0.605 against 0.738.
\item[Inference] The errors are partly decorrelated, and Branch~A's complementary strength is
concentrated in Zoom\_voice, a minority class, which is consistent with the recall
offsets of Experiment~7.
\item[Verdict] Blending is worth testing; weight selection is done nested.
\end{explog}

\begin{explog}{23}{Blend weight: full-data sweep and nested grouped selection}
\item[Hypothesis] A weighted average $w\,p_\mathrm{A}+(1-w)\,p_\mathrm{B}$ has an interior
optimum, and if the weight is real, the one chosen on the training side of each
grouped fold should be stable and its held-out score should hold.
\item[Method] $w$ was swept from 0.00 to 1.00 in steps of 0.05 on all 1{,}285 rows, with
and without the Zoom rule (reference only: weight and score share the same rows).
Then, for each of the ten grouped folds, $w$ was chosen from the same grid using only
that fold's training rows and applied unchanged to its held-out rows, under two
selection objectives (macro-F1 without the rule, and macro-F1 after the rule).
Held-out predictions are pooled over both repeats (each row scored twice).
\item[Result] Full-data sweep: without the rule best $w=0.40$, macro-F1 0.8363 ($w=0$:
0.8201; $w=0.5$: 0.8351; $w=1$: 0.7900); with the rule best $w=0.50$, 0.8380
($w=0.40$: 0.8366; $w=0$: 0.8292; $w=1$: 0.7855). Nested selection:
Table~\ref{tab:t3-nested}. The rule-free objective selects 0.45, 0.30, 0.40, 0.40,
0.45, 0.40, 0.45, 0.40, 0.25, 0.45 (median 0.40); the rule-aware objective selects
0.55, 0.25, 0.50, 0.55, 0.30, 0.55, 0.50, 0.40, 0.50, 0.50 (median 0.50).
\item[Inference] The optimum is broad (0.40--0.50), the weight is stable across folds
(0.25--0.55) and the held-out gain over Branch~B alone survives nesting. The
full-data figures are optimistic and are used only to locate the region.
\item[Verdict] Weight 0.40/0.60 adopted (the nested median under the rule-free objective
and the full-data optimum of the sweep). The headline estimate is the per-fold mean
of the nested blend with the rule, \textbf{0.8318}. This figure inherits the stratified-fold optimism of Branch~A's probabilities (Section~\ref{sec:t3-limits}).
\end{explog}

\begin{table}[t]
\caption{Nested grouped selection of the blend weight (Experiment~23). Pooled: both
repeats of every held-out row; fold mean: mean of the ten per-fold scores.}
% TODO: confirm that the fold means 0.8286 / 0.8318 belong to the rule-free objective (as in the final script header)
\label{tab:t3-nested}
\centering
\begin{tabular}{lcccc}
\hline
 & \multicolumn{2}{c}{Macro-F1 pooled} & \multicolumn{2}{c}{Macro-F1 fold mean} \\
Configuration & no rule & + rule & no rule & + rule \\
\hline
Branch B alone (TabICL, flat+hier) & 0.8201 & 0.8292 & 0.8175 & 0.8268 \\
Blend, nested weight, rule-free objective (median 0.40) & 0.8319 & 0.8345 & 0.8286 & \textbf{0.8318} \\
Blend, nested weight, rule-aware objective (median 0.50) & 0.8307 & 0.8331 & --- & --- \\
Blend, fixed $w=0.40$ (optimistic, same rows) & 0.8363 & 0.8366 & --- & --- \\
\hline
\end{tabular}
\end{table}
 
\begin{explog}{24}{The rule after the blend, and the blend against Branch B}
\item[Hypothesis] The rule is orthogonal to which model produced the probabilities, so it
should add to the blend as it does to Branch~B; and the blend improves on Branch~B as
shipped (with the rule).
\item[Method] Rule applied to the blend, to Branch~B and to Branch~A alone, same predictions
otherwise. Paired bootstrap over rows (2{,}000 resamples, both repeats averaged) and
McNemar's test on row correctness, shipped configuration ($w=0.40$ plus rule)
against Branch~B plus rule and against raw Branch~B.
\item[Result] Blend: 0.8363 $\rightarrow$ 0.8366 (fixed $w=0.40$) and 0.8319 $\rightarrow$
0.8345 (nested), i.e.\ $+0.0003$ to $+0.0026$. Branch~B alone: 0.8201 $\rightarrow$
0.8292 ($+0.0091$). Branch~A alone: 0.7900 $\rightarrow$ 0.7855 ($-0.0045$). Against
Branch~B plus rule: $+0.0072$ macro-F1 (pooled 0.8366 against 0.8292; interval
$[-0.0036,+0.0192]$, $P(\Delta>0)=0.888$); McNemar 16 rows won by the blend against
10 for Branch~B, $p=0.327$. Against raw Branch~B: $+0.0164$ (interval
$[-0.0007,+0.0330]$, $P(\Delta>0)=0.970$). With the nested weight estimate the gain
over Branch~B plus rule is $+0.0039$ (pooled) and $+0.0050$ (per-fold mean).
\item[Inference] The rule's value is largest where the model under-predicts Zoom\_voice;
Branch~A already assigns Zoom\_voice at high recall, so most of the rule's effect is
absorbed by the blend, and applied to Branch~A alone it costs Zoom\_video recall. The
blend is better in every estimate, by roughly 0.4 to 0.7 macro-F1 points over
Branch~B with the rule, but the interval includes zero; most of the improvement over
\emph{raw} Branch~B is the Zoom rule plus Branch~A's Zoom\_voice recall.
\item[Verdict] Blend and rule both retained (threshold 300 bytes, no fitted parameter);
the marginal gain of each over the best single branch is reported as not established.
\end{explog}

\begin{explog}{25}{Final configuration and test-set output}
\item[Hypothesis] The validated pipeline can be refit on all training rows and applied to the 327
test flows without change.
\item[Method] Branch~A (31 features; the three tree models; $T=2$; prior correction) was
refit on all 1{,}285 rows, single-threaded; the per-class offsets and the specialist
of Experiment~7 were \emph{not} carried into the blend, so that the blend receives
corrected probabilities rather than outputs shifted toward recall. Branch~B flat and
hierarchical test probabilities (16 estimators, full-data refit) were averaged, with
the label order aligned to Branch~A's alphabetical order. Blend $0.4\,p_\mathrm{A}+0.6\,p_\mathrm{B}$,
argmax, then the Zoom rule.
\item[Result] The rule routed 43 of 327 test flows to Zoom\_voice, in line with the training
proportion (an even 80/80 split in the ambiguous subset). Predicted class counts:
Discord\_voice 67, Discord\_video 56, Zoom\_voice 43, Messenger\_video 33,
Zoom\_video 27, GoogleMeet\_video 26, Messenger\_voice 23, WhatsApp\_video 22,
WhatsApp\_voice 19, GoogleMeet\_voice 11 (these match the submitted file exactly).
\item[Inference] The output proportions are plausible against the flows-per-call prior of
Experiment~1.
\item[Verdict] Submitted. Nested-CV estimate: macro-F1 0.8318.
\end{explog}


\subsection{Per-Class Results}

Table~\ref{tab:t3-perclass} gives the per-class breakdown of the fixed-weight blend.

\begin{table}[t]
\caption{Per-class results of the fixed-weight configuration ($w=0.40$ plus Zoom
rule; pooled, so the mean F1 is 0.8366, not the nested 0.8318) against Branch~B plus
rule and raw Branch~A.}
\label{tab:t3-perclass}
\centering
\begin{tabular}{lrcccc}
\hline
Class & $n$ & Precision & Recall & F1 & F1 B+rule / A raw \\
\hline
Discord\_video & 256 & 0.969 & 0.855 & 0.909 & 0.905 / 0.846 \\
Discord\_voice & 238 & 0.837 & 0.903 & 0.869 & 0.864 / 0.826 \\
GoogleMeet\_video & 112 & 0.814 & 0.741 & 0.776 & 0.765 / 0.737 \\
GoogleMeet\_voice & 40 & 0.720 & 0.900 & 0.800 & 0.759 / 0.679 \\
Messenger\_video & 131 & 0.929 & 0.893 & 0.911 & 0.914 / 0.867 \\
Messenger\_voice & 94 & 0.907 & 0.936 & 0.921 & 0.906 / 0.873 \\
WhatsApp\_video & 82 & 0.888 & 0.963 & 0.924 & 0.936 / 0.865 \\
WhatsApp\_voice & 80 & 0.988 & 1.000 & 0.994 & 0.988 / 0.975 \\
Zoom\_video & 172 & 0.978 & 0.512 & 0.672 & 0.672 / 0.703 \\
Zoom\_voice & 80 & 0.437 & 0.912 & 0.591 & 0.583 / 0.530 \\
\hline
\end{tabular}
\end{table}

\subsection{Limitations}
\label{sec:t3-limits}
 
\begin{itemize}
\item \textbf{Mixed fold protocols.} Branch~A's out-of-fold probabilities were
produced under stratified 5-fold cross-validation, whereas Branch~B's come from the
grouped folds, so Branch~A's raw scores in Experiments~22--24 are likely somewhat
optimistic relative to a grouped estimate. Correcting this would weaken Branch~A, so
the selected weight of 0.40 on Branch~A is, if anything, too high, and the nested
blend estimate of 0.8318 is likely somewhat optimistic.
% TODO [RERUN]: recompute Branch A on the stored grouped folds and redo Tables 10 and 11
\item \textbf{Same-row selection.} The Branch~B variant, the 16-estimator setting and
the blend were all chosen among near-equal candidates on the same 1{,}285 rows; every
paired difference in Experiments~16 and 24 lies within bootstrap noise.
\item \textbf{Different tuned configurations.} Branch~A's tree settings were tuned
for the 31-column set under its own protocol and differ from the Experiment~3 search
winners, which were found on the 241-column set.
\item \textbf{Row-level resampling.} Bootstrap intervals resample rows, not
pseudo-groups, and are therefore slightly narrow.
\item \textbf{Scope of the rule.} The Zoom rule is a property of this dataset's
Zoom audio-only windows (Experiment~18); it is not claimed to transfer to other
applications, where voice and video are separable (AUC 0.83--0.88).
\item \textbf{External and unlabelled data.} Experiment~12 used a public capture
corpus (MIRAGE-AppAct-2024) for pretraining and probes only; the final pipeline does
not use it. Unlabelled test flows were used for adversarial validation and the
domain-adversarial pretraining head, never for fitting a label-dependent quantity.
\end{itemize}
